\chapter{Hipotez Testleri, Hata Türleri ve Etki Büyüklüğü}
\label{ch:hypothesis}
\index{hipotez testi}
\index{etki büyüklüğü}

\begin{hedefler}
Bu bölümün sonunda okuyucu:
\begin{itemize}
  \item araştırma sorusunu hedef parametre üzerinden sıfır ve alternatif hipoteze dönüştürebilecek,
  \item test istatistiğini ve $p$-değerini doğru biçimde açıklayabilecek,
  \item anlamlılık düzeyi, istatistiksel karar ve bilimsel yorum arasındaki farkı görebilecek,
  \item Tip I ve Tip II hatayı PDR bağlamında yorumlayıp test gücüyle ilişkilendirebilecek,
  \item tek ve çift yönlü test seçimini araştırma planına göre gerekçelendirebilecek,
  \item $p$-değeri, güven aralığı ve etki büyüklüğünü birlikte değerlendirebilecek,
  \item çoklu test, seçici raporlama ve veri görüldükten sonra hipotez değiştirme risklerini tanıyabilecek,
  \item hipotez testi sonuçlarını Python ve SPSS çıktılarından okuyup şeffaf biçimde raporlayabilecektir.
\end{itemize}
\end{hedefler}

\begin{hazirlik}
``$p=0{,}03$ ise sıfır hipotezinin doğru olma olasılığı yüzde 3'tür''
ifadesindeki kavram yanılgısını belirleyin. Bu olasılık ifadesinde koşulun
hangi yönde kurulması gerektiğini düşünün.
\end{hazirlik}

\section{PDR vakası: Danışmanlık programı işe yarıyor mu?}

Bir üniversite, birinci sınıf öğrencilerinin yalnızlık düzeyini azaltmak için
altı haftalık grup danışmanlığı programı uygulamıştır. Program grubunun
son-test yalnızlık puanı ortalaması karşılaştırma grubundan 4,2 puan daha
düşüktür. ``Program işe yaradı'' sonucuna yalnız bu gözlenen farktan hareketle
ulaşılabilir mi?

Örneklem seçimi yeniden yapılsaydı grup ortalamaları ve aralarındaki fark bir
miktar değişirdi. Hipotez testi, gözlenen farkın sıfır fark varsayımı altında
ne kadar olağandışı olduğunu değerlendirir. Bununla birlikte araştırmacının
yanıtlaması gereken sorular bundan daha geniştir:

\begin{enumerate}
  \item Karşılaştırılan hedef parametre tam olarak nedir?
  \item Gözlenen fark rastlantısal örnekleme değişkenliğiyle açıklanabilir mi?
  \item Farkın büyüklüğü öğrenciler açısından anlamlı mıdır?
  \item Araştırma tasarımı nedensel bir yorum yapmaya izin vermekte midir?
  \item Sonuç hangi hata ve yanlılık kaynaklarına duyarlıdır?
\end{enumerate}

Bu bölümde test işlemi bir ``anlamlılık üretme'' aracı olarak değil, önceden
tanımlanmış bir araştırma sorusunu belirsizlik altında değerlendirme yöntemi
olarak ele alınacaktır.

\begin{researchbox}
İstatistiksel olarak belirgin bir azalma, programın her öğrenciye yarar
sağladığını veya programın tek başına değişime neden olduğunu göstermez.
Nedensel yorum için karşılaştırma grubu, rastgele atama, uygulama bütünlüğü,
kayıp katılımcılar ve eş zamanlı müdahaleler de değerlendirilmelidir.
\end{researchbox}

\section{Araştırma sorusundan sınanabilir hipoteze}
\index{sıfır hipotezi}
\index{alternatif hipotez}

Hipotezler örneklemdeki gözlenen sayılar hakkında değil, evren parametreleri
hakkında kurulur. Program ve karşılaştırma gruplarının evren ortalamaları
sırasıyla $\mu_P$ ve $\mu_K$ olsun. Fark ``program eksi karşılaştırma'' olarak
tanımlandığında çift yönlü hipotezler
\[
H_0:\mu_P-\mu_K=0,
\qquad
H_1:\mu_P-\mu_K\neq0
\]
biçimindedir. Programın yalnızlığı azalttığı yönü önceden ve güçlü biçimde
gerekçelendirilmişse alternatif hipotez $H_1:\mu_P-\mu_K<0$ olarak tek yönlü
kurulabilir.

\subsection{Sıfır ve alternatif hipotezin rolleri}

\begin{description}
  \item[Sıfır hipotezi] Test istatistiğinin referans dağılımını kurmak için kullanılan kesin iddiadır; çoğunlukla farkın ya da ilişkinin olmadığını belirtir.
  \item[Alternatif hipotez] Verinin destekleyip desteklemediği karşıt değerler kümesidir; çift yönlü veya önceden belirlenmiş tek yönlü olabilir.
\end{description}

Sıfır hipotezinin her zaman ``hiçbir şey olmaz'' biçiminde kurulması gerekmez.
Bir ölçeğin ortalamasının norm değeri 50 olup olmadığı, bir program etkisinin
en az 3 puana ulaşıp ulaşmadığı veya iki yöntemin uygulamada eşdeğer olup
olmadığı farklı sıfır değerleri ve farklı test tasarımları gerektirebilir.

\begin{table}[htbp]
\centering
\caption{Araştırma sorusundan hipoteze geçiş örnekleri.}
\begin{tabular}{@{}p{0.30\textwidth}p{0.23\textwidth}p{0.28\textwidth}@{}}
\toprule
Araştırma sorusu & Hedef parametre & Olası hipotez \\
\midrule
Program grupları farklı mı? & $\mu_P-\mu_K$ & $H_0:\mu_P-\mu_K=0$ \\
Başvuru oranı yüzde 20'den yüksek mi? & $p$ & $H_0:p=0{,}20$, $H_1:p>0{,}20$ \\
Kaygı ile uyku ilişkili mi? & $\rho$ & $H_0:\rho=0$, $H_1:\rho\neq0$ \\
İki kategorik değişken bağımsız mı? & Dağılımsal ilişki & $H_0$: değişkenler bağımsızdır \\
\bottomrule
\end{tabular}
\end{table}

Hipotezler veri görülmeden önce yazılmalıdır. Örneklem sonucu 4,2 puanlık
azalma gösterdikten sonra ``zaten azalma bekliyorduk'' diyerek çift yönlü
testi tek yönlüye çevirmek yanlış pozitif riskini artırır.

\section{Hipotez testinin mantığı}
\index{test istatistiği}
\index{anlamlılık düzeyi}

Çoğu klasik test aynı düşünceyi izler:

\begin{enumerate}
  \item Hedef parametreyi, $H_0$ ve $H_1$ hipotezlerini tanımla.
  \item Anlamlılık düzeyini ve analiz yöntemini önceden belirle.
  \item Varsayımları ve veri kalitesini incele.
  \item Gözlenen tahmini ve standart hatayı hesapla.
  \item Test istatistiğini $H_0$ altındaki referans dağılımla karşılaştır.
  \item $p$-değeri ve güven aralığıyla istatistiksel kanıtı değerlendir.
  \item Etki büyüklüğü, uygulamalı önem ve tasarım sınırlılıklarıyla bilimsel yorumu tamamla.
\end{enumerate}

Birçok test istatistiğinin genel yapısı
\[
\text{test istatistiği}
=\frac{\text{gözlenen tahmin}-\text{$H_0$ değeri}}
{\text{tahminin standart hatası}}
\]
biçimindedir. Mutlak değer büyüdükçe gözlenen sonuç $H_0$ değerinden standart
hata birimleriyle daha uzakta bulunur.

\begin{sezgibox}
Test istatistiği bir ``sinyal/gürültü'' oranı gibi düşünülebilir. Paydaki fark
sinyali, paydadaki standart hata ise örnekleme gürültüsünü temsil eder. Aynı
fark daha küçük standart hatayla gözlendiğinde $H_0$ altında daha olağandışı
görünür.
\end{sezgibox}

\section{Test istatistiği ve p-değeri}
\index{p-değeri}
\index{test istatistiği}

$p$-değeri, $H_0$ ve test modelinin diğer varsayımları doğru kabul edildiğinde,
gözlenen test istatistiği kadar veya ondan daha uç bir sonuç elde etme
olasılığıdır. Çift yönlü testte ``daha uç'' her iki yöndeki sonuçları kapsar
\cite{casella2002,wasserman2004}.

Program örneğinde ortalama fark $-4{,}2$ ve farkın standart hatası 1,7 olsun.
Sıfır fark hipotezi altında
\[
t=\frac{-4{,}2-0}{1{,}7}\approx-2{,}47
\]
elde edilir. Serbestlik derecesine bağlı kesin değer küçük ölçüde değişmekle
birlikte çift yönlü $p$-değeri yaklaşık 0,015 civarındadır. Bu sonuç, gerçek
evren farkı sıfırsa gözlenen kadar veya daha uç bir test istatistiğinin seyrek
olacağını gösterir.

\begin{figure}[htbp]
\centering
\begin{tikzpicture}[alt={Çift yönlü hipotez testinde sıfır dağılımının iki kuyruğunda p-değerini oluşturan alanlar},x=1cm,y=1cm]
  \draw[->,PrismGray] (-4.3,0)--(4.5,0) node[right,font=\scriptsize]{test istatistiği};
  \fill[PrismWarnOrange!65,smooth,samples=60,domain=2.47:4]
    (2.47,0)--plot (\x,{2.5*exp(-\x*\x/2)})--(4,0)--cycle;
  \fill[PrismWarnOrange!65,smooth,samples=60,domain=-4:-2.47]
    (-4,0)--plot (\x,{2.5*exp(-\x*\x/2)})--(-2.47,0)--cycle;
  \draw[PrismBlue,very thick,smooth,samples=100,domain=-4:4]
    plot (\x,{2.5*exp(-\x*\x/2)});
  \draw[PrismWarnOrange!90!black,dashed] (-2.47,0)--(-2.47,0.13);
  \draw[PrismWarnOrange!90!black,dashed] (2.47,0)--(2.47,0.13);
  \node[font=\scriptsize,text=PrismGray] at (0,1.2) {$H_0$ altındaki dağılım};
  \node[font=\scriptsize,text=PrismWarnOrange!80!black] at (3.25,0.42) {$p/2$};
  \node[font=\scriptsize,text=PrismWarnOrange!80!black] at (-3.25,0.42) {$p/2$};
\end{tikzpicture}
\caption{Çift yönlü testte $p$-değerini oluşturan iki uç alan.}
\end{figure}

\subsection{$p$-değeri ne değildir?}

$p$-değeri:

\begin{itemize}
  \item $H_0$'ın doğru olma olasılığı değildir,
  \item sonucun yalnız şanstan kaynaklanma olasılığı değildir,
  \item etkinin büyüklüğünü göstermez,
  \item sonucun yeniden üretileceği olasılığı vermez,
  \item araştırma tasarımındaki yanlılıkları veya varsayım ihlallerini ölçmez.
\end{itemize}

``$p=0{,}03$'' ifadesi, $H_0$ doğru kabul edildiğinde bu kadar veya daha uç
verinin olasılığının yüzde 3 olduğunu söyler. Aranan $P(H_0\mid\text{veri})$
değil, hesaplanan $P(\text{veri kadar uç sonuç}\mid H_0)$ olasılığıdır. Bu iki
koşullu olasılık birbirinin yerine kullanılamaz.

\begin{mistakebox}
$p=0{,}049$ ile $p=0{,}051$ bilimsel açıdan iki zıt dünya oluşturmaz. Bu
değerler verinin $H_0$ ile uyumu bakımından birbirine çok yakındır. Keskin
``var/yok'' sınıflaması yerine kesin $p$-değeri, güven aralığı ve etki
büyüklüğü birlikte incelenmelidir.
\end{mistakebox}

\section{Anlamlılık düzeyi ve karar dili}
\index{anlamlılık düzeyi}
\index{istatistiksel karar}

Anlamlılık düzeyi $\alpha$, $H_0$ doğru olduğunda uzun dönemde göze alınan Tip
I hata oranıdır. Yaygın $\alpha=0{,}05$ seçimi doğal bir yasa değildir;
yanlış pozitif ve yanlış negatif kararların sonuçlarına göre araştırma
planında gerekçelendirilmelidir.

Önceden belirlenen karar kuralında $p\leq\alpha$ ise $H_0$ reddedilir;
$p>\alpha$ ise $H_0$'ı reddetmek için yeterli kanıt bulunmadığı söylenir.
İkinci durumda ``$H_0$ doğrudur'' veya ``gruplar aynıdır'' denmez.

İstatistiksel karar ile bilimsel sonuç ayrılmalıdır:

\begin{description}
  \item[İstatistiksel karar] Önceden tanımlı test kuralının sonucudur: reddet veya reddetme.
  \item[Kanıtın gücü] $p$-değeri, güven aralığının konumu ve genişliğiyle dereceli olarak değerlendirilir.
  \item[Bilimsel yorum] Etki büyüklüğü, tasarım, ölçme niteliği, önceki bilgi ve uygulama bağlamını içerir.
\end{description}

\section{Güven aralığı ile test arasındaki ilişki}
\index{güven aralığı!hipotez testi ilişkisi}

Aynı model ve standart hata yöntemi kullanıldığında, yüzde 95 güven aralığının
sıfır farkını dışlaması $\alpha=0{,}05$ düzeyindeki çift yönlü testte
$H_0$'ın reddedilmesine karşılık gelir. Program farkı için önceki bölümde
hesaplanan
\[
[-7{,}53;\ -0{,}87]
\]
aralığı sıfırı dışlar ve yaklaşık $p=0{,}015$ sonucuyla tutarlıdır.

Güven aralığı test kararından daha fazla bilgi verir. Yalnız ``sıfır
dışlandı'' demek yerine verinin yaklaşık 0,9 ile 7,5 puanlık azalmalarla
uyumlu olduğu görülür. Aralığın genişliği, etkinin büyüklüğü hakkındaki
belirsizliği de ortaya koyar.

\begin{etkinlik}
Üç çalışma için fark güven aralıkları $[-8;-4]$, $[-5;2]$ ve $[-0{,}5;0{,}4]$
olsun. Her birini (a) sıfır farkına ilişkin kanıt, (b) tahmin hassasiyeti ve
(c) uygulamada önemli olabilecek değerler bakımından karşılaştırın.
\end{etkinlik}

\section{Tip I ve Tip II hata}
\index{Tip I hata}
\index{Tip II hata}
\index{alfa@$\alpha$}
\index{beta@$\beta$}

Test kararı kesin bilgiye değil, örnekleme dayandığı için iki tür hata
oluşabilir:

\begin{table}[htbp]
\centering
\caption{Test kararları ve olası hata türleri.}
\begin{tabular}{@{}p{0.25\textwidth}p{0.27\textwidth}p{0.27\textwidth}@{}}
\toprule
 & $H_0$ doğru & $H_0$ yanlış \\
\midrule
$H_0$ reddedilir & Tip I hata olasılığı $\alpha$ & Doğru karar; olasılığı güç \\
$H_0$ reddedilmez & Doğru karar & Tip II hata olasılığı $\beta$ \\
\bottomrule
\end{tabular}
\end{table}

\subsection{PDR bağlamında hata sonuçları}

Program gerçekte etkisizken etkili sonucuna varmak Tip I hatadır. Kurum bu
durumda zaman ve bütçeyi yararsız bir programa ayırabilir. Program gerçekte
yararlı olduğu hâlde yeterli kanıt bulamamak Tip II hatadır; öğrenciler etkili
bir hizmetten mahrum kalabilir.

Hataların sonuçları bağlama göre eşit değildir. Zararlı yan etkisi veya yüksek
maliyeti bulunan bir müdahalede yanlış olumlu karar daha ağır görülebilir.
Erken tarama araştırmasında ise önemli bir riski kaçırmak daha ağır olabilir.
Bu değerlendirme $\alpha$, hedef güç ve örneklem hacmi planına yansıtılmalıdır.

\begin{mistakebox}
$\alpha=0{,}05$, yürütülen belirli çalışmada yanlış karar verme olasılığının
kesin olarak yüzde 5 olduğu anlamına gelmez. Bu değer, $H_0$ doğru olduğunda
aynı karar kuralının uzun dönemli yanlış ret oranını tanımlar.
\end{mistakebox}

\section{Test gücü}
\index{test gücü}
\index{güç analizi}

Test gücü $1-\beta$ olup belirli bir gerçek etki var olduğunda $H_0$'ı
reddetme olasılığıdır. Güç tek bir sabit sayı değildir; hangi büyüklükteki
gerçek etkinin saptanmak istendiğine bağlıdır \cite{cohen1988}.

Gücü artıran başlıca etkenler şunlardır:

\begin{itemize}
  \item daha büyük gerçek etki,
  \item daha büyük ve dengeli örneklem,
  \item daha düşük ölçüm ve bireysel değişkenlik,
  \item güvenilir ölçme aracı,
  \item bağımlı tasarımda yüksek ve uygun eşleşme,
  \item diğer koşullar aynıyken daha yüksek $\alpha$ veya önceden gerekçeli tek yönlü test.
\end{itemize}

Son iki seçenek yanlış pozitif riskiyle ödünleşim içerir; gücü artırmanın en
sağlıklı yolu çoğunlukla yeterli örneklem, kaliteli ölçüm ve güçlü araştırma
tasarımıdır. Araştırma başlamadan önce uygulamada anlamlı en küçük etki,
beklenen değişkenlik, hedef güç ve anlamlılık düzeyi kullanılarak güç analizi
yapılabilir. Sosyal bilimlerde yüzde 80 güç sık kullanılan bir başlangıç
hedefidir; bağlama göre daha yüksek bir değer gerekebilir.

\begin{sezgibox}
Düşük güçlü bir test, loş ışıkta küçük bir nesneyi aramaya benzer. Nesneyi
görememek onun bulunmadığını kanıtlamaz. Geniş güven aralığı da çalışmanın hem
önemsiz hem önemli etkilerle uyumlu olabileceğini gösterir.
\end{sezgibox}

\section{Tek ve çift yönlü test}
\index{tek yönlü test}
\index{çift yönlü test}

Çift yönlü test $H_1:\theta\neq\theta_0$ biçiminde her iki yöndeki sapmayı
değerlendirir. Tek yönlü test $H_1:\theta>\theta_0$ veya
$H_1:\theta<\theta_0$ biçiminde yalnız önceden belirtilen yönü sınar.

Tek yönlü test ancak şu koşullarda savunulabilir:

\begin{itemize}
  \item yön kuramsal ve uygulamalı olarak veri görülmeden önce belirlenmiştir,
  \item ters yöndeki sonuç araştırma açısından gerçekten aynı iddiayı desteklemeyecektir,
  \item test yönü analiz planında açıkça kaydedilmiştir.
\end{itemize}

Program yalnızlığı artırırsa bu sonuç da bilimsel ve etik açıdan önemliyse
çift yönlü test daha uygundur. Veri görüldükten sonra daha küçük $p$-değeri
elde etmek amacıyla test yönü değiştirilemez.

\section{Etki büyüklüğü}
\index{etki büyüklüğü}

Etki büyüklüğü farkın veya ilişkinin örneklem büyüklüğünden daha bağımsız bir
ölçekte ne kadar büyük olduğunu gösterir. $p$-değeri ``veri sıfır etkiyle ne
kadar uyumlu?'' sorusuna, etki büyüklüğü ise ``gözlenen fark veya ilişki ne
kadar büyük?'' sorusuna yaklaşır.

\subsection{Cohen'in $d$ değeri}
\index{Cohen d@Cohen'in $d$ değeri}

Bağımsız iki grup ortalaması için standartlaştırılmış fark
\[
d=\frac{\bar x_1-\bar x_2}{s_{\text{birleşik}}}
\]
biçimindedir. Program farkı $-4{,}2$ ve birleşik standart sapma 10 ise
\[
d=-0{,}42
\]
elde edilir. İşaret yönü, mutlak değer standartlaştırılmış büyüklüğü gösterir.

Cohen'in yaklaşık 0,2, 0,5 ve 0,8 eşikleri bazen küçük, orta ve büyük olarak
kullanılır; ancak bunlar evrensel karar sınırları değildir. Ölçeğin niteliği,
müdahalenin maliyeti, önceki çalışmalar ve öğrenciler açısından anlamlı en
küçük değişim önceliklidir.

\subsection{Diğer yaygın ölçüler}

\begin{table}[htbp]
\centering
\caption{Bazı araştırma soruları için yaygın etki büyüklükleri.}
\begin{tabular}{@{}p{0.27\textwidth}p{0.19\textwidth}p{0.34\textwidth}@{}}
\toprule
Araştırma yapısı & Ölçü & Temel yorum \\
\midrule
İki ortalamanın farkı & Cohen'in $d$ değeri & Standart sapma biriminde fark \\
İki nicel değişken & Pearson $r$ & Doğrusal ilişkinin yönü ve derecesi \\
ANOVA & $\eta^2$, $\omega^2$ & Sonuç değişkenliğinin grupla ilişkili payı \\
Kategorik ilişki & Cramér $V$ & Çapraz tablodaki ilişkinin derecesi \\
İki oran veya risk & Risk farkı, risk oranı & Mutlak veya göreli olasılık değişimi \\
\bottomrule
\end{tabular}
\end{table}

Etki büyüklüğü de örneklemden tahmin edilir ve belirsizdir. Mümkün olduğunda
güven aralığıyla birlikte raporlanmalıdır. Küçük bir örneklemde $d=0{,}50$
bulunması, gerçek etkinin tam 0,50 olduğu anlamına gelmez.

\begin{researchbox}
PDR uygulamasında ortalama etki bireysel farklılıkları gizleyebilir. Ortalama
azalma orta büyüklükte olsa bile bazı öğrenciler güçlü yarar, bazıları değişim
ve bazıları kötüleşme yaşayabilir. Etki büyüklüğü klinik değişim, zarar,
katılım ve eşitlik göstergeleriyle birlikte değerlendirilmelidir.
\end{researchbox}

\section{İstatistiksel anlamlılık ve uygulamalı önem}
\index{istatistiksel anlamlılık}
\index{uygulamalı önem}

Çok büyük örneklemde önemsiz küçük etkiler düşük $p$-değeri üretebilir. Küçük
örneklemde ise uygulamada önemli bir etki geniş güven aralığı nedeniyle
istatistiksel olarak belirgin bulunmayabilir. Bu nedenle sonuç dört bileşenle
birlikte okunmalıdır:

\begin{enumerate}
  \item nokta tahmini ve yönü,
  \item etki büyüklüğü,
  \item güven aralığının konumu ve genişliği,
  \item $p$-değeri ve araştırma tasarımı.
\end{enumerate}

Program için uygulamada anlamlı en küçük azalmanın 3 puan olduğu varsayılsın.
$[-7{,}53;-0{,}87]$ aralığı hem 3 puandan büyük hem de 3 puandan küçük
azalmaları içerir. Sonuç sıfır farkına karşı kanıt sunsa da uygulamada önemli
etkinin büyüklüğü konusunda hâlâ belirsizlik vardır.

\section{Çoklu test sorunu}
\index{çoklu test}
\index{aile düzeyinde hata oranı}
\index{yanlış keşif oranı}
\index{Bonferroni düzeltmesi}
\index{Holm düzeltmesi}

Her biri $\alpha=0{,}05$ düzeyinde 20 bağımsız sıfır hipotezi sınanırsa, en az
bir yanlış pozitif bulma olasılığı yaklaşık
\[
1-(1-0{,}05)^{20}\approx0{,}64
\]
olur. Testler bağımsız olmasa da temel sorun değişmez: çok sayıda fırsat,
tesadüfen düşük bir $p$-değeri bulma olasılığını artırır. Aile düzeyindeki
hata ve yanlış keşif oranı bu çokluk sorununa farklı hedeflerle yaklaşır
\cite{tukey1953,benjamini1995}.

Sorunu yönetmek için:

\begin{itemize}
  \item birincil ve ikincil sonuçlar önceden belirlenebilir,
  \item test ailesi açıkça tanımlanabilir,
  \item Bonferroni veya daha güçlü Holm düzeltmesi kullanılabilir,
  \item keşfedici analizler doğrulayıcı analizlerden ayrılabilir,
  \item bütün sonuçlar yalnız anlamlı olanlar seçilmeden raporlanabilir.
\end{itemize}

Bonferroni yönteminde $m$ test için her test yaklaşık $\alpha/m$ düzeyinde
değerlendirilir. Yöntem basit fakat özellikle ilişkili çok sayıda testte
koruyucu olabilir. Düzeltme seçimi analiz amacına ve yanlış keşif maliyetine
göre yapılmalıdır.

\section{Seçici analiz ve araştırmacı esnekliği}
\index{seçici raporlama}
\index{araştırmacı esnekliği}
\index{ön kayıt}

Aynı veri üzerinde farklı dışlama kuralları, sonuç değişkenleri, alt gruplar,
kovaryatlar ve test yönleri denenip yalnız düşük $p$-değerinin raporlanması
yanlış pozitif oranını görünmez biçimde yükseltir. Bu uygulamalar bazen
``$p$-değeri avcılığı'' olarak adlandırılır.

Şeffaflığı artırmak için araştırmacı:

\begin{itemize}
  \item hipotezleri ve temel analiz planını veri incelemesinden önce kaydetmeli,
  \item veri toplamanın ne zaman duracağını önceden belirlemeli,
  \item sonuçlara göre yapılan değişiklikleri açıkça belirtmeli,
  \item duyarlılık analizlerini ana analizden ayırmalı,
  \item anlamlı ve anlamsız bütün ilgili sonuçları raporlamalıdır.
\end{itemize}

Ön kayıt kötü bir tasarımı iyi hâle getirmez ve keşfedici araştırmayı
yasaklamaz. Amacı, önceden öngörülen doğrulayıcı sınamalar ile veriden doğan
yeni hipotezlerin birbirinden ayırt edilmesini sağlamaktır.

\begin{mistakebox}
Veriden ilham alan bir hipotez bilimsel olarak değersiz değildir; fakat aynı
veri hem hipotezi üretmek hem de onu bağımsız biçimde doğrulamak için
kullanılmış gibi sunulmamalıdır. Yeni hipotez sonraki veriyle sınanmalıdır.
\end{mistakebox}

\section{Test varsayımları ve sağlamlık}

Her testin sonucu kullanılan veri üretim modeli altında geçerlidir. Yaygın
kontroller şunlardır:

\begin{itemize}
  \item örneklem seçimi ve gruplara atama süreci,
  \item gözlemlerin bağımsızlığı veya bağımlılığın modellenmesi,
  \item değişkenin ölçme düzeyi ve puanlama niteliği,
  \item aykırı değerler, dağılım biçimi ve varyans yapısı,
  \item eksik verinin miktarı ve oluşma mekanizması,
  \item analiz planıyla veri yapısının uyumu.
\end{itemize}

Varsayım kontrolü, sonuç anlamlı değilse alternatif test arama işlemi
değildir. Veri yapısına uygun yöntem araştırma sorusu ve tasarımla birlikte
seçilmelidir. Parametrik olmayan bir yöntem de otomatik olarak varsayımsız
değildir; yalnız farklı varsayımlara ve hedeflere dayanır.

\begin{assumptionbox}
Öğrenciler sınıflar içinde kümelenmişse sıradan bağımsızlık varsayımı standart
hataları küçültebilir. Çok küçük $p$-değeri bile yanlış standart hata veya
yanlı tasarımı telafi etmez. Önce analiz birimini ve bağımlılık yapısını
doğrulayın.
\end{assumptionbox}

\section{Python ile bütünleştirici analiz}

Aşağıdaki örnek iki bağımsız grubun yalnızlık puanlarını Welch $t$ testiyle
karşılaştırır, ortalama farkı ve Cohen'in $d$ değerini hesaplar:

\begin{lstlisting}
import numpy as np
from scipy import stats

program = np.array([38, 41, 35, 44, 39, 36, 40, 42,
                    33, 37, 45, 39, 34, 41, 38])
karsilastirma = np.array([43, 47, 39, 48, 42, 45, 41, 46,
                         44, 40, 49, 43, 47, 42, 46])

sonuc = stats.ttest_ind(program, karsilastirma,
                        equal_var=False)
fark = program.mean() - karsilastirma.mean()

n1, n2 = len(program), len(karsilastirma)
s1, s2 = program.std(ddof=1), karsilastirma.std(ddof=1)
birlesik_s = np.sqrt(((n1 - 1) * s1**2 +
                      (n2 - 1) * s2**2) /
                     (n1 + n2 - 2))
d = fark / birlesik_s

print("Ortalamalar:", program.mean(), karsilastirma.mean())
print("Ortalama farki:", fark)
print("Welch t testi:", sonuc)
print("Cohen d:", d)
\end{lstlisting}

Kod çalıştırıldığında yalnız $p$-değeri raporlanmamalıdır. Grup ortalamaları
ve standart sapmaları, farkın yönü, serbestlik derecesi, güven aralığı ve etki
büyüklüğü de verilmelidir. Gerçek araştırmada veri temizleme ve varsayım
incelemesi test komutundan önce gelir.

\section{SPSS ile hipotez testi çıktısını okuma}

Bağımsız iki grup ortalaması için:

\begin{enumerate}
  \item \textbf{Analyze $>$ Compare Means $>$ Independent-Samples T Test} yolunu izleyin.
  \item Sonuç değişkenini \textbf{Test Variable(s)}, grup değişkenini \textbf{Grouping Variable} alanına taşıyın.
  \item \textbf{Define Groups} ile grup kodlarını doğru tanımlayın.
  \item Önce grup $n$, ortalama ve standart sapmalarını kontrol edin.
  \item Varyans yapısına uygun satırdan $t$, serbestlik derecesi, iki yönlü $p$, ortalama farkı ve güven aralığını okuyun.
  \item Farkın SPSS tarafından hangi grup eksi hangi grup olarak hesaplandığını doğrulayın.
\end{enumerate}

Levene testinin anlamlı olmaması varyansların kesin eşit olduğunu kanıtlamaz.
Welch satırı, özellikle grup büyüklükleri veya varyanslar farklı olduğunda
sağlam bir varsayılan seçenek olabilir. Yöntem seçimi yalnız Levene
$p$-değerine bağlı mekanik bir karar hâline getirilmemelidir.

\begin{outputguidebox}
Çıktıyı şu sırayla okuyun: (1) grup tanımları ve örneklem sayıları, (2)
betimsel istatistikler, (3) farkın yönü, (4) test istatistiği ve serbestlik
derecesi, (5) kesin $p$-değeri, (6) farkın güven aralığı, (7) ayrıca
hesaplanan etki büyüklüğü. ``Sig.'' sütununu tek başına sonuç olarak
raporlamayın.
\end{outputguidebox}

\section{Bütünleştirici karar örneği}

Program eksi karşılaştırma farkı $-4{,}2$, standart hata 1,7, test
istatistiği yaklaşık $-2{,}47$, çift yönlü $p=0{,}015$, yüzde 95 güven aralığı
$[-7{,}53;-0{,}87]$ ve Cohen'in $d=-0{,}42$ olsun.

\begin{enumerate}
  \item $p<0{,}05$ olduğu ve aralık sıfırı dışladığı için sıfır farkına karşı istatistiksel kanıt vardır.
  \item Negatif işaret program grubunun ortalamasının daha düşük olduğunu gösterir.
  \item $d=-0{,}42$ küçük ile orta arasında standartlaştırılmış bir farka işaret eder; bağlamsal yorum gerekir.
  \item Güven aralığı etkinin küçük veya uygulamada önemli olabilecek büyüklüklerde olabileceğini gösterir.
  \item Nedensel sonuç, ancak atama ve izleme tasarımı bunu destekliyorsa kurulabilir.
\end{enumerate}

Bu değerlendirme ``program anlamlıdır'' cümlesinden daha bilgilendiricidir.
İstatistiksel yöntem, etkiyi, belirsizliği ve tasarımın izin verdiği yorum
sınırını birlikte görünür kılmalıdır.

\section{Bilimsel raporlama örneği}

``Program grubunun son-test yalnızlık puanı karşılaştırma grubundan ortalama
4,2 puan daha düşüktür. Welch bağımsız örneklem $t$ testi gruplar arasında
istatistiksel olarak belirgin bir fark göstermiştir, $t(sd)=-2{,}47$,
$p=0{,}015$, yüzde 95 GA $[-7{,}53;-0{,}87]$. Standartlaştırılmış fark
$d=-0{,}42$'dir. Güven aralığı küçük ve uygulamada önemli olabilecek
azalmaları birlikte içerdiğinden etkinin büyüklüğü ihtiyatla yorumlanmıştır.''

Gerçek raporda $sd$ yerine yazılımın verdiği serbestlik derecesi, her grubun
$n$, ortalama ve standart sapması yazılmalıdır. $p$-değeri çok küçükse
$p<0{,}001$ biçimi kullanılabilir; diğer durumlarda mümkün olduğunca kesin
değer verilmelidir.

\section{Bir hipotez testini değerlendirirken}

\begin{enumerate}
  \item Hipotezler hedef parametre üzerinden ve veri görülmeden önce kurulmuş mu?
  \item Test araştırma tasarımına, ölçme düzeyine ve bağımlılık yapısına uygun mu?
  \item Anlamlılık düzeyi ve tek--çift yön seçimi önceden belirlenmiş mi?
  \item Tahmin, güven aralığı ve etki büyüklüğü $p$-değeriyle birlikte verilmiş mi?
  \item Aralık uygulamada önemli değerleri içeriyor mu?
  \item Çoklu testler ve analiz planından sapmalar açıklanmış mı?
  \item İstatistiksel sonuç nedensellik veya klinik karar sınırlarını aşmadan yorumlanmış mı?
\end{enumerate}

\section{Literatür veri setiyle IBM SPSS laboratuvarı}
\index{Student Performance veri seti!hipotez testi}
\index{SPSS!tek örneklem t testi@tek örneklem $t$ testi}
\index{SPSS!binom testi}
\index{etki büyüklüğü!Cohen d@Cohen $d$}
\index{etki büyüklüğü!Cohen h@Cohen $h$}

Bu uygulamada \emph{Student Performance} matematik dosyasındaki yıl sonu notu
\texttt{G3} ve bu nottan üretilen başarı göstergesi kullanılacaktır
\cite{cortez2008data,cortezsilva2008}. Aynı veri iki farklı hedef parametre
üzerinden sınanacaktır:

\begin{enumerate}
  \item Evren ortalama G3 puanı 10'dan farklı mıdır?
  \item G3 puanı en az 10 olan öğrencilerin evren oranı yüzde 50'den farklı mıdır?
\end{enumerate}

İki sınama da veri incelenmeden önce belirlenmiş çift yönlü hipotezler ve
$\alpha=0{,}05$ anlamlılık düzeyiyle yürütülecektir. Referans değerler yalnız
öğretim amacıyla kullanılmaktadır; 10 puanının veya yüzde 50 oranının resmî
bir politika hedefi olduğu ileri sürülmemektedir.

\begin{researchbox}
Ortalama not ile başarı oranı farklı hedef parametrelerdir. Aynı veri
dosyasından hesaplanmaları aynı test sonucunu vermelerini gerektirmez. Başarı
oranı ayrıca önceden seçilen 10 puanlık eşiğe bağlıdır; eşik değiştiğinde oran
testi değişirken ham not ortalaması değişmez.
\end{researchbox}

\subsection{Hipotezler ve analiz planı}

Ortalama için
\[
H_0:\mu_{G3}=10,
\qquad
H_1:\mu_{G3}\neq10
\]
hipotezleri tek örneklem $t$ testiyle; başarı oranı için
\[
H_0:p_{\text{başarı}}=0{,}50,
\qquad
H_1:p_{\text{başarı}}\neq0{,}50
\]
hipotezleri tek örneklem binom testiyle değerlendirilecektir. Ortalama için
Cohen'in $d$ değeri, oran için referans orandan fark ve Cohen'in $h$ değeri
etki büyüklüğü olarak raporlanacaktır.

\subsection{SPSS syntax}

Başarı göstergesi eksik G3 değerlerini yanlışlıkla başarısız olarak
kodlamamak için önce sistem eksik olarak oluşturulur. Tek örneklem $t$ testi
ortalama farkını, binom testi ise gözlenen başarı sayısını yüzde 50 referans
oranıyla karşılaştırır.

\begin{lstlisting}[language={}]
DATASET ACTIVATE StudentMath.

COMPUTE basarili=$SYSMIS.
IF NOT MISSING(G3) basarili=(G3>=10).
VARIABLE LABELS basarili 'Yıl sonu notu en az 10'.
VALUE LABELS basarili 0 'Başarısız' 1 'Başarılı'.
VARIABLE LEVEL basarili (NOMINAL).
FORMATS basarili (F1.0).
EXECUTE.

T-TEST
 /TESTVAL=10
 /MISSING=ANALYSIS
 /VARIABLES=G3
 /CRITERIA=CI(.95).

NPAR TESTS
 /BINOMIAL (.50)=basarili
 /MISSING ANALYSIS.
\end{lstlisting}

Etki büyüklükleri aşağıdaki syntax ile doğrudan veri özetlerinden
hesaplanabilir. \texttt{MODE=ADDVARIABLES} özetleri veri dosyasındaki bütün
satırlara eklediği için \texttt{LIST} yalnız ilk satırı gösterir.

\begin{lstlisting}[language={}]
AGGREGATE
 /OUTFILE=* MODE=ADDVARIABLES
 /BREAK=
 /n_G3=N(G3)
 /ortalama_G3=MEAN(G3)
 /ss_G3=SD(G3)
 /n_basarili=N(basarili)
 /x_basarili=SUM(basarili).

COMPUTE fark_ortalama=ortalama_G3-10.
COMPUTE cohen_d=fark_ortalama/ss_G3.
COMPUTE p_hat=x_basarili/n_basarili.
COMPUTE fark_oran=p_hat-.50.
COMPUTE cohen_h=2*ARSIN(SQRT(p_hat))
 -2*ARSIN(SQRT(.50)).
FORMATS ortalama_G3 ss_G3 fark_ortalama cohen_d
 p_hat fark_oran cohen_h (F7.3).

LIST VARIABLES=n_G3 ortalama_G3 ss_G3
 fark_ortalama cohen_d n_basarili x_basarili
 p_hat fark_oran cohen_h
 /CASES=FROM 1 TO 1.
\end{lstlisting}

Menüyle ortalama testi için \textbf{Analyze $>$ Compare Means $>$ One-Sample
T Test} yolunda G3 test değişkeni, 10 ise \textbf{Test Value} olarak girilir.
Binom testi için SPSS sürümüne göre \textbf{Analyze $>$ Nonparametric Tests $>$
Legacy Dialogs $>$ Binomial} yolu kullanılabilir; test oranı 0,50 olarak
belirlenir. Menü işlemleri \textbf{Paste} düğmesiyle syntax'a aktarılıp analiz
günlüğünde saklanmalıdır.

\subsection{SPSS çıktısının erişilebilir özeti}

\begin{table}[htbp]
\centering
\caption{G3 ortalaması ve başarı oranı için hipotez testi sonuçları.}
\small
\begin{tabular}{@{}p{0.27\textwidth}rrrrrr@{}}
\toprule
Sınama & $n$ & Tahmin & Fark & İstatistik & $p$ & Etki \\
\midrule
G3 ortalaması & 395 & 10,415 & 0,415 & $t(394)=1{,}801$ & 0,072 & $d=0{,}091$ \\
Başarı oranı & 395 & 0,671 & 0,171 & $x=265$ & $<0{,}001$ & $h=0{,}349$ \\
\bottomrule
\end{tabular}
\end{table}

Ortalama testinde SPSS'nin \textbf{Mean Difference} değeri
$10{,}415-10=0{,}415$'tir. Farkın yüzde 95 güven aralığı yaklaşık
$[-0{,}038;0{,}868]$ olup sıfırı içerir. Başarı oranında ise 395 öğrencinin
265'i başarı eşiğini karşılamış; gözlenen oran yüzde 67,1 olmuştur. Binom
testinin çift yönlü kesin $p$-değeri çok küçük olduğundan tabloda
$p<0{,}001$ biçiminde raporlanmıştır.

\subsection{Çıktının analizi}

G3 ortalaması referans değer 10'dan 0,415 puan yüksektir. Bununla birlikte
$t(394)=1{,}801$ ve $p=0{,}072$ olduğundan yüzde 5 anlamlılık düzeyinde
$H_0$'a karşı yeterli istatistiksel kanıt yoktur. Bu karar ortalamanın tam
olarak 10 olduğunu veya iki değer arasında hiç fark bulunmadığını kanıtlamaz.
Güven aralığı, yaklaşık 0,04 puanlık azalma ile 0,87 puanlık artış arasındaki
farkların veriyle uyumlu olduğunu gösterir.

Ortalama için $d=0{,}091$ değeri, gözlenen farkın yaklaşık onda bir standart
sapma olduğunu ve örneklemde çok küçük bir standartlaştırılmış farka işaret
ettiğini gösterir. Büyük örneklem bulunmasına rağmen $p$-değerinin 0,05'in
altına düşmemesi yalnız örneklem hacmiyle değil, gözlenen etkinin küçük
olmasıyla da ilişkilidir.

Başarı oranı yüzde 50 referansından 17,1 yüzde puan yüksektir. Binom testi
$H_0:p=0{,}50$ hipotezine karşı güçlü istatistiksel kanıt verir
($p<0{,}001$). Önceki bölümde hesaplanan Wilson yüzde 95 güven aralığı
$[0{,}623;0{,}715]$ olup 0,50'yi dışlar ve test kararıyla uyumludur.
Cohen'in $h=0{,}349$ değeri farkın yalnız istatistiksel değil, dikkate değer
bir büyüklükte olabileceğini gösterir; ancak eğitimsel önemi bağlam ve eşik
tanımıyla ayrıca değerlendirilmelidir.

Sonuçların farklı görünmesi bir çelişki değildir. Ortalama testi bütün not
büyüklüklerini ve 10 çevresindeki farkı kullanırken binom testi her notu
yalnızca eşik altında veya üstünde olarak sınıflandırır. Bu nedenle araştırma
sorusu ve hedef parametre test seçilmeden önce açıkça belirlenmelidir.

\begin{mistakebox}
$p=0{,}072$ sonucu ``ortalama kesinlikle 10'dur'', $p<0{,}001$ sonucu da
``başarı oranı kesinlikle yüzde 67,1'dir'' anlamına gelmez. İlk sonuç sıfır
hipotezine karşı kanıtın seçilen eşikte yetersiz olduğunu, ikinci sonuç ise
yüzde 50 referansıyla güçlü bir uyumsuzluk bulunduğunu gösterir. Parametre
tahmini için güven aralıkları ayrıca okunmalıdır.
\end{mistakebox}

\begin{outputguidebox}
Çıktıyı şu sırayla okuyun: (1) hipotez ve referans değer, (2) geçerli $n$,
(3) gözlenen tahmin ve farkın yönü, (4) test istatistiği, (5) kesin
$p$-değeri, (6) güven aralığı, (7) etki büyüklüğü, (8) tasarım ve eşik
sınırlılıkları. Anlamlı ve anlamsız sonuçları yalnız 0,05'in iki tarafında
yer aldıkları için niteliksel olarak bütünüyle farklıymış gibi sunmayın.
\end{outputguidebox}

\subsection{Örnek bulgu paragrafı}

``İki okuldan gözlenen 395 öğrencinin yıl sonu matematik notu ortalaması,
10 referans değerinden istatistiksel olarak belirgin biçimde farklı değildir
($\bar x=10{,}42$, $SS=4{,}58$), $t(394)=1{,}80$, $p=0{,}072$, ortalama
farkı için yüzde 95 GA $[-0{,}04;0{,}87]$, $d=0{,}09$. Buna karşılık notu en
az 10 olan 265 öğrencinin oranı yüzde 67,1'dir ve yüzde 50 referans oranından
farklıdır, çift yönlü kesin binom $p<0{,}001$, Wilson yüzde 95 GA
$[\%62{,}3;\%71{,}5]$, $h=0{,}35$. İki analiz farklı hedef parametreleri
sınadığından sonuçlar birbiriyle çelişkili yorumlanmamıştır. Örneklemin iki
okulla sınırlı olması nedeniyle ülke geneline doğrudan genelleme
yapılmamıştır.''

\section{R ile hipotez testi}
\label{sec:r-b09}

\textbf{Soru.} Bölümdeki çift yönlü testte gözlenen ortalama farkı 2,1,
standart hata 1 ve serbestlik derecesi 49 olsun. $H_0:\delta=0$ için
\pval-değerini, yüzde 95 güven aralığını ve $\alpha=0{,}05$ kararını bulun.

\begin{lstlisting}[style=prismR]
fark <- 2.1
standart_hata <- 1
serbestlik <- 49
alfa <- 0.05
t_degeri <- fark / standart_hata
p_degeri <- 2 * pt(abs(t_degeri), df = serbestlik,
                   lower.tail = FALSE)
kritik <- qt(1 - alfa / 2, df = serbestlik)
aralik <- fark + c(-1, 1) * kritik * standart_hata
c(t = t_degeri, p = p_degeri, kritik = kritik)
aralik
if (p_degeri < alfa) {
  print("H0 reddedilir")
} else {
  print("H0 reddedilemez")
}
pt(t_degeri, df = serbestlik, lower.tail = FALSE)
\end{lstlisting}

\begin{outputguidebox}
\textbf{Çözüm ve kontrol değerleri.} $t=2{,}10$, çift yönlü
$p=0{,}04090$ ve kritik değer 2,0096'dır. Güven aralığı
$[0{,}0904;4{,}1096]$ sıfırı içermez; yüzde 5 düzeyinde $H_0$ reddedilir.
Son satır, önceden belirlenmiş $H_1:\delta>0$ için tek yönlü
$p=0{,}02045$ verir.
\end{outputguidebox}

\textbf{Yorum.} Çift yönlü değer, mutlak test istatistiğinin üst kuyruk
olasılığının iki katıdır. Sonuca bakıp yön seçmek uygun değildir. Küçük
\pval-değeri $H_0$'ın doğru olma olasılığı veya etkinin önemli olduğu anlamına gelmez.

\textbf{Pekiştirme ve yanıt.} Önceden $\alpha=0{,}01$ seçilseydi aynı
$p=0{,}04090$ ile $H_0$ reddedilemezdi. Veri değişmez; karar eşiği değişir.
\section{R ile test gücü ve tek örneklem testi}
\label{sec:r-b10}

\textbf{Soru.} $n=25$, $\bar x=55$, $s=10$ için referans ortalama 50'yi
çift yönlü test edin. Ayrıca veri toplanmadan önce gerçek farkın 5, evren
standart sapmasının 10 olduğu varsayımıyla $n=34$ için gücü hesaplayın.

\begin{lstlisting}[style=prismR]
hacim <- 25
ortalama <- 55
standart_sapma <- 10
referans <- 50
standart_hata <- standart_sapma / sqrt(hacim)
t_degeri <- (ortalama - referans) / standart_hata
p_degeri <- 2 * pt(abs(t_degeri), df = hacim - 1,
                   lower.tail = FALSE)
aralik <- ortalama + c(-1, 1) *
  qt(0.975, df = hacim - 1) * standart_hata
c(t = t_degeri, p = p_degeri,
  d = (ortalama - referans) / standart_sapma)
aralik
power.t.test(n = 34, delta = 5, sd = 10,
             sig.level = 0.05, type = "one.sample",
             alternative = "two.sided", strict = TRUE)
plan <- power.t.test(delta = 5, sd = 10, power = 0.80,
                     sig.level = 0.05, type = "one.sample",
                     alternative = "two.sided", strict = TRUE)
ceiling(plan$n)
\end{lstlisting}

\begin{outputguidebox}
\textbf{Çözüm ve kontrol değerleri.} $t(24)=2{,}50$, $p=0{,}01965$,
$d=0{,}50$ ve yüzde 95 güven aralığı $[50{,}8722;59{,}1278]$'dir.
Planlama varsayımları altında 34 gözlemin gücü yaklaşık 0,8078'dir.
Yüzde 80 hedef güç için hesaplanan hacim yukarı yuvarlandığında 34 bulunur.
\end{outputguidebox}

\textbf{Yorum.} \texttt{strict = TRUE} çift yönlü testin iki ret bölgesini
de güç hesabına katar. Güç hesabındaki 5 puan ve 10 puan önceden gerekçelenmiş
planlama değerleri olmalıdır; gözlenen anlamlılığı yeniden ifade eden
``gözlenen güç'' olarak sunulmamalıdır.

\textbf{Pekiştirme ve yanıt.} Tip II hata olasılığı bu belirli alternatif
ve tasarım altında yaklaşık $1-0{,}8078=0{,}1922$'dir. Daha küçük bir gerçek
fark için aynı örneklem hacminin gücü daha düşük olur.
\section{Bölüm sonu çözümlü problemler}

\begin{enumerate}
  \item \textbf{Test istatistiği ve karar.} Gözlenen fark 4,5, sıfır hipotezi değeri 0 ve standart hata 1,5'tir. Test istatistiğini bulun; çift yönlü $p=0{,}006$ ise sonucu yazın.

  \textbf{Çözüm:} $t=(4{,}5-0)/1{,}5=3$ olur. $p=0{,}006<0{,}05$ olduğundan sıfır farkına karşı kanıt vardır. Sonuç $H_0$'ın doğru olma olasılığının yüzde 0,6 olduğu anlamına gelmez.

  \item \textbf{Hata türleri.} Gerçekte etkisiz bir programı etkili ilan etmek ve gerçekte etkili bir programı etkisiz saymak hangi hata türleridir?

  \textbf{Çözüm:} Etkisiz programı etkili ilan etmek yanlış pozitif, yani Tip I hatadır. Etkili program için yeterli kanıt bulamamak yanlış negatif, yani Tip II hatadır. Tip II hata olasılığı $\beta$, güç ise $1-\beta$'dır.

  \item \textbf{Çoklu test.} Bağımsız olduğu varsayılan 10 testin her biri $\alpha=0{,}05$ ile yürütülürse en az bir yanlış pozitif bulma olasılığını yaklaşık hesaplayın.

  \textbf{Çözüm:} $1-(1-0{,}05)^{10}=1-0{,}95^{10}\approx0{,}401$'dir. Olasılık yaklaşık yüzde 40,1'e yükselir. Testler bağımlı olduğunda kesin değer değişir; çoklu karşılaştırma sorunu yine sürer.
\end{enumerate}

\bolumkaynaklari{\cite{casella2002,cohen1988,benjamini1995,tukey1953,cortez2008data,cortezsilva2008}}

\section*{Hatırla--Uygula--Yansıt}

\begin{enumerate}
\renewcommand{\labelenumi}{\textbf{\Alph{enumi}.}}
  \item \textbf{Hatırla:} $H_0$, $H_1$, $\alpha$, $p$-değeri, $\beta$ ve güç kavramlarını birbirinden ayırın.
  \item \textbf{Hipotez kur:} Bir PDR programının kaygı puanını değiştirip değiştirmediği için hedef parametreyi ve çift yönlü hipotezleri yazın.
  \item \textbf{Hesapla:} Gözlenen fark 6, $H_0$ değeri 0 ve standart hata 2 ise test istatistiğini bulun.
  \item \textbf{Yorumla:} $p=0{,}03$ sonucunun ne söylediğini ve söylemediğini iki ayrı cümleyle açıklayın.
  \item \textbf{Karar dili:} $p=0{,}18$ için ``$H_0$ kabul edilir'' yerine bilimsel olarak daha uygun bir ifade yazın.
  \item \textbf{Hata türleri:} Etkisiz bir programı uygulamaya alma ve etkili bir programı reddetme durumlarını Tip I--Tip II hata ile eşleştirin.
  \item \textbf{Güç:} Etki büyüklüğü, örneklem hacmi, değişkenlik ve $\alpha$ değiştiğinde gücün nasıl etkileneceğini açıklayın.
  \item \textbf{Yön seç:} Programın hem yarar hem zarar olasılığının önemli olduğu bir araştırmada neden çift yönlü test seçileceğini tartışın.
  \item \textbf{Etki büyüklüğü:} Ortalama fark 5 ve birleşik standart sapma 12 ise Cohen'in $d$ değerini hesaplayıp bağlama bağlı yorumlayın.
  \item \textbf{Aralıkla yorumla:} Fark için yüzde 95 güven aralığı $[-1;9]$ olduğunda ``etki yoktur'' sonucunun neden yetersiz olduğunu açıklayın.
  \item \textbf{Çoklu test:} 15 sonuç değişkenini ayrı ayrı sınamanın yanlış pozitif riskini ve iki olası çözümü açıklayın.
  \item \textbf{Uygula:} Python kodundaki veriyi değiştirerek örneklem hacmi ve grup farkının $p$-değeri ile $d$ üzerindeki etkisini inceleyin.
  \item \textbf{Eleştir:} Yalnız ``$p<0{,}05$ bulundu'' cümlesiyle sunulan bir araştırma sonucunda eksik olan bilgileri listeleyin.
  \item \textbf{Yansıt:} İstatistiksel olarak anlamlı bir ortalama değişimin neden her öğrenci için klinik veya eğitsel yarar anlamına gelmediğini tartışın.
  \item \textbf{Raporla:} Tahmin, güven aralığı, test istatistiği, serbestlik derecesi, $p$-değeri, etki büyüklüğü ve yorum sınırını içeren kısa bir bulgu paragrafı yazın.
\end{enumerate}